% Part: first-order-logic
% Chapter: introduction
% Section: semantic-notions

\documentclass[../../../include/open-logic-section]{subfiles}

\begin{document}

\olfileid{fol}{int}{sem}

\olsection{అర్థపర భావనలు}

$\Sat{M}{!A}[s]$ నిలుస్తుందా అని పరిశీలించేటప్పుడు, సౌలభ్యం కోసం $s$
అన్ని \tetoken{చరాలకు}{!!{variable}s} విలువలు నిర్దేశిస్తుందని తీసుకుంటాం; కానీ
$!A$లోని \tetoken{చరాలకు}{!!{variable}s} అది నిర్దేశించే విలువలను మాత్రమే వాడతామని
పైన చెప్పాం. నిజానికి, $!A$లోని \emph{స్వేచ్ఛా} చరరాశుల విలువలు
మాత్రమే ముఖ్యం. మనం జాగ్రత్తగా వ్యవహరిస్తాం కాబట్టి, ఈ వాస్తవాన్ని
నిరూపిస్తాం. \tetoken{వాక్యాలకు}{!!{sentence}s} స్వేచ్ఛా చరరాశులు ఉండవు కనుక, అవి ఒక
\tetoken{నిర్మాణంలో}{!!a{structure}} సంతృప్తి చెందుతాయా లేదా అన్న విషయంలో $s$కు అసలు
ప్రాధాన్యం లేదు. కాబట్టి $!A$ ఒక \tetoken{వాక్యం}{!!a{sentence}} అయినప్పుడు,
``అన్ని~$s$లకు $\Sat{M}{!A}[s]$'' అనే అర్థంలో $\Sat{M}{!A}$ను
నిర్వచించవచ్చు; అది కనీసం ఒక~$s$కు $\Sat{M}{!A}[s]$ కావడానికీ
సరిసమానమని తేలుతుంది. \tetoken{సూత్రాలకు}{!!{formula}s} పనిచేసే సంతృప్తి నిర్వచనం
పొందడానికి \tetoken{చరం}{!!{variable}} నిర్దేశాలను ప్రవేశపెట్టాలి; కానీ
\tetoken{వాక్యాలకు}{!!{sentence}s} సంతృప్తి, \tetoken{చరం}{!!{variable}} నిర్దేశాలపై ఆధారపడదు.

``$\Sat{M}{!A}$''కు నిర్వచనం వచ్చిన తరువాత, మొదటిస్థాయి విధేయ తర్కపు
\tetoken{వాక్యాలు}{!!{sentence}s} విషయంలో ``సందర్భం,'' ``దానిలో సత్యం'' అంటే ఏమిటో
తెలుస్తుంది. \tetoken{వాక్యాలు}{!!{sentence}s} కోసం $\Sat{M}{!A}$ నిర్వచనం ఆధారంగా
చెల్లుబాటుతనం, అనుగమనం, సంతృప్తిపరచదగినతనం అనే ప్రాథమిక అర్థపర
భావనలను నిర్వచించవచ్చు. ప్రతి \tetoken{నిర్మాణం}{!!{structure}} ఒక వాక్యాన్ని
సంతృప్తిపరిస్తే అది చెల్లుబాటైనది, $\Entails !A$. ఒక
\tetoken{వాక్యాలు}{!!{sentence}s} సమితి~$\Gamma$లోని అన్ని \tetoken{వాక్యాలను}{!!{sentence}s}
సంతృప్తిపరిచే ప్రతి \tetoken{నిర్మాణం}{!!{structure}}, $!A$ను కూడా సంతృప్తిపరిస్తే,
ఆ సమితి నుంచి అది అనుగమిస్తుంది, $\Gamma \Entails !A$. ఏదో ఒక
\tetoken{నిర్మాణం}{!!{structure}}, ఒక \tetoken{వాక్యాలు}{!!{sentence}s} సమితిలోని అన్ని \tetoken{వాక్యాలను}{!!{sentence}s}
ఒకేసారి సంతృప్తిపరిస్తే ఆ సమితి సంతృప్తిపరచదగినది.

\tetoken{సూత్రాలు}{!!{formula}s} ఆగమనాత్మకంగా నిర్వచించబడ్డాయి; సంతృప్తి కూడా
\tetoken{సూత్రాలు}{!!{formula}s} నిర్మాణంపై ఆగమనం ద్వారా నిర్వచించబడుతుంది. అందువల్ల మన
అర్థవిచార ధర్మాలను నిరూపించడానికి, నిర్వచించిన అర్థపర భావనలను
పరస్పరం అనుసంధానించడానికి ఆగమనాన్ని ఉపయోగించవచ్చు. ఈ ధర్మాల్లో
కొన్నింటిని సమీకరించి నిరూపిస్తాం. అవి ఒక్కొక్కటీ ఆసక్తికరమైనవి
కావడం కొంత కారణం; కానీ అర్థవిచారానికి, \tetoken{వ్యుత్పత్తి}{!!{derivation}} వ్యవస్థలకు
మధ్య సంబంధాన్ని తరువాత పరిశీలించినప్పుడు వాటిలో అనేకం ఉపయోగపడతాయనేది
ప్రధాన కారణం. అలా చేయడానికి, మనం ఇంకా ప్రస్తావించని కొన్ని
వాక్యనిర్మాణ భావనలు, క్రియలను కూడా (ఖచ్చితంగా, అంటే ఆగమనం ద్వారా)
నిర్వచించాలి.

\end{document}
